\section{Preliminaries and normalization}
\label{sec:normalization}

The constructions below operate on finite, simple, undirected input graphs.
This entails no restriction on the theorem if the input convention permits
loops or parallel edge identities. Replace each input edge by a path of
length three with two fresh internal vertices. For a common edge use the
same path in both colours; for a private edge give the new vertices
isolated copies in the other colour. The resulting graphs are simple and
have total size $O(N)$. A SEFE extends by choosing subdivision points away
from all other drawing features. Conversely, suppress the path's internal
vertices in its colour or colours and delete their isolated copies in the
other colour. This restores the original edge curves and is an exact
linear-size reduction. Empty graphs and union components with at most two
vertices after this step are handled directly.

Auxiliary SPQR skeletons and Synchronized Planarity instances may be loopless
multigraphs; a digon is then a simple cycle. Drawings are considered on an
oriented sphere. Selecting a point outside the finite drawing converts the
result to the plane. A block is a maximal biconnected subgraph or a bridge.
An isolated vertex belongs to no edge block.

A rotation system specifies a cyclic order of incident darts at every
vertex. It is planar if each connected component has Euler characteristic
$|V|-|E|+f=2$, where faces are obtained by tracing the dart successor
permutation, with $f=1$ for an isolated-vertex component. For a connected graph this determines its oriented spherical
embedding. For a disconnected graph one must additionally record the face
of each component containing every other component. These data can be
represented by a component--region containment tree. Equal rotations alone
are therefore insufficient in the disconnected case.

We use the common-embedding criterion: a pair has a SEFE exactly when
planar embeddings of its two colours induce the same embedding of $H$,
including relative component positions \cite[Theorem~4]{junger}. To see the sufficiency,
identify orientation-preserving regular neighbourhoods of corresponding
embedded common components. Their matching face positions identify the
remaining complementary regions. Extend the boundary identifications
through these planar regions to an orientation-preserving homeomorphism.
Applying it to one colour identifies the whole common drawing. Isolated
common vertices are matched as marked points in the corresponding regions.
The standard topological facts behind this sketch are stated at the end of
Section~\ref{rep:global}.

For a cyclic order $\sigma$ and distinct elements $a,d,b$, put
\[
 \orcyc_\sigma(a,d,b)=
 \begin{cases}0,&d\text{ occurs between }a\text{ and }b\text{ going forward},\\
 1,&\text{otherwise}.
 \end{cases}
\]
Reflection complements this bit. A directed cycle $C$ has left side $0$
and right side $1$. With outgoing and incoming darts $c_{\rm out},c_{\rm in}$
at $v$, the local gap of a noncycle dart $d$ is
$\gap(C,v,d)=\orcyc_{\rho_v}(c_{\rm out},d,c_{\rm in})$.
The connected component entered by $d$ after deleting $C$ is on that side.
For three leaves in a plane tree, their frontier orientation equals the
orientation of their three directions at their median. This follows by
reading the three disjoint contour excursions from the median.

\subsection{The published infrastructure used}

We use the classical block-cut and SPQR decompositions and their exact
embedding parameterization: R-skeleton reflections and P-skeleton cyclic
orders, opposite at their two poles, determine all block embeddings.
Substitution along twin virtual edges is the constructive inverse. We also
use ordinary deterministic planarity testing. These established algorithms
are sufficient at any polynomial running time; their linear bounds are
used only to simplify the estimates.

The other imported infrastructure is the ordinary Synchronized Planarity
algorithm of Bl\"asius--Fink--Rutter (BFR) \cite{bfr}: its graph operations,
their legal applicability, its reduction schedule and potential decrease.
Section~\ref{cse:section} proves the additional marking and observer
invariants, forward and backward, rather than assuming that the ordinary
algorithm preserves those extra constraints. The ordinary target has
Q-cells synchronizing fixed rotations up to reflection, and pipes requiring
full reversed correspondences between P-vertex rotations. The pipes form
a matching. No unrestricted simultaneous-ordering oracle is used.

\subsection{Connecting the colours without changing old common blocks}

First solve the connected components of $U=G_1\cup G_2$ independently.
Their solutions can occupy disjoint disks, so this is an equivalence in
both directions. In a connected union, if colour $i$ is disconnected,
some private edge $v_1v_2$ of the other colour joins distinct components
of $G_i$. The following fact is the connectivity normalization of
Bl\"asius--Rutter \cite[Section~2, Lemma~1 and Theorem~1]{brdisconnect}; we
give the nesting argument because it matters for arbitrary disconnected $H$.

\begin{lemma}[Cofacial relocation]
\label{norm:cofacial}
If an instance has a SEFE and an opposite-colour private edge $v_1v_2$
joins different components of $G_i$, it has a SEFE in which $v_1,v_2$
are incident with a common face of $G_i$.
\end{lemma}
\begin{proof}
In a SEFE the interior of the other-colour edge is in one face $f$ of
$H$. Near its endpoints it lies in faces $g_1,g_2$ of $G_i$, both
contained in $f$. If $g_1=g_2$, nothing is required. Otherwise let $Q_j$
be the colour-$i$ component containing $v_j$, and let $b$ be the face of
$Q_2$ containing $Q_1$. Let $\mathcal Y$ consist of $Q_2$ and all
colour-$i$ components lying in the complement of $b$.

Choose a small closed disk $\Delta'$ in $g_2$ near $v_2$, and a small
closed disk $\Delta$ in $g_1$ near $v_1$. Choose an orientation-preserving
sphere homeomorphism taking the complement of the interior of $\Delta'$
onto $\Delta$. Apply it to exactly the components in $\mathcal Y$.
They all fit in the empty disk $\Delta$, while all other colour-$i$
components remain outside. This is a planar drawing with the same
component rotations. The region outside $\Delta$ in $g_1$ and the image
of $g_2\setminus\operatorname{int}\Delta'$ join across $\partial\Delta$;
hence $v_1$ and the moved $v_2$ are cofacial.

It remains to verify the common embedding. Consider common components
$X$ outside $\mathcal Y$ and $X'$ inside. Before the move,
the complement of $b$, together with a small neighbourhood at $v_2$,
is connected, avoids $X$, contains $X'$, and meets the private edge
in $f$. Thus $X'$ is in the face of $X$ containing $f$. Afterwards it
lies in $\Delta\subset f$, so this position is unchanged.
Conversely, either $b$ itself, if $X'\subseteq Q_2$, or
$b\cup\Gamma_i(Q_2)$ otherwise, connects $X$ to $g_1$ while avoiding
$X'$. Its position relative to $X'$ is consequently the face containing
$f$. After the move, all unchanged components lie in the image of
$\operatorname{int}\Delta'$, the same transported face. Pairs both inside
or both outside retain their positions. The common rotations also
remain unchanged. The common-embedding criterion completes the SEFE.
\end{proof}

Add a fresh vertex $z$, a common edge $v_1z$, and a private colour-$i$
edge $zv_2$. This is an exact augmentation. For the forward direction use
Lemma~\ref{norm:cofacial}, put $z$ near $v_1$ in the shared local sector,
draw the small common segment $v_1z$ there, and complete the colour-$i$
arc to $v_2$ in its cofacial region. The other colour only receives the
small pendant common segment, adjacent to its existing edge $v_1v_2$.
For the local-sector claim, choose a disk at $v_1$ disjoint from the moved
components and from the relocation disk $\Delta$. The common drawing is
unchanged there. The common-embedding homeomorphism can be chosen to be
the identity on this disk: extend its complementary-region boundary maps
relative to the additional fixed disk. Thus the old opposite-colour edge
germ and its colour-$i$ sector $g_1$ are unchanged, and the two local
sectors used for the new common segment agree.
The backward direction deletes the new vertex and its edges.

The number of colour-$i$ components decreases by one; that of the other
colour does not increase. At most $2(|V|-1)$ steps suffice, so the total
size remains $O(N)$. Each step adds only a common pendant bridge. Every
old nontrivial common block, and hence its SPQR type, is unchanged.
No bridge-to-triangle replacement is required.

\subsection{Splitting the union and preserving all relative positions}

Suppose both colours are connected and $w$ is a union cutvertex. For each
$D\in\cc(U-w)$, restrict both colours to $D\cup\{w\}$. Each restriction
is connected: a simple colour path from a vertex of $D$ to $w$ cannot
leave $D$ before it reaches $w$. Restriction of a SEFE gives solutions
of all these factors.

Conversely, take compatible embeddings of all factors. At $w$, cut both
rotations of a factor immediately after the same common dart, if one
exists; otherwise cut each arbitrarily. Concatenate the cut rotations
in the same order in both colours. Factor dart sets are intervals, so
the resulting colour rotations are planar. Common rotations agree.

Every common cycle $C$ lies in one factor. Vertices in another factor
lie on the side of $w$ if $w\notin C$. If $w\in C$, they lie in the
chosen corner after that factor's common dart; the side of this corner
depends only on the common rotation and is the same in both colours.
Within the cycle's own factor the sides already agree. Distinct faces
of a connected common component are distinguished by the faces of its
blocks containing them (Lemma~\ref{rep:lem:face-signatures}, which applies
to every connected plane graph), and distinct faces of a biconnected block
are distinguished by their sides with respect to a cycle basis of that
block (Lemma~\ref{gen:basis}). Hence they are distinguished by the
block-cycle side information. Therefore all relative common-component
positions agree after gluing. This proves the converse for arbitrary
$H$, not only for a cactus.

Repeated splitting yields union-block factors of total linear size,
counting repeated cutvertices. A common biconnected block lies wholly
in one union block, so its SPQR structure is preserved. A union bridge
must be common once both colours are connected; its factor is trivial.
The remaining factors have biconnected union and connected colours.
Thus this normalization is an exact, class-preserving reduction for
Theorem~\ref{thm:main}.

\subsection{NP membership of unrestricted SEFE-2}

The same normalization works without the restriction on common blocks.
For every resulting factor, a certificate consists of the two colour
rotation systems. These have $O(N)$ darts and $O(N\log(N+1))$ encoded
bits. Face tracing verifies colour planarity. Restricting the rotations
checks equality of common orders.

For each nontrivial common block $K$, search the components of $G_i-V(K)$.
Their faces relative to $K$ can be read at an attachment: in the rotation
at that vertex, the sector containing the attachment dart identifies a
face walk of the restricted embedding of $K$. Connectedness and colour
planarity place the whole searched component in that face. Compare the
two colours' face labels for every vertex outside $K$. Bridges have one
face. The argument in Section~\ref{rep:section} shows that equal common
rotations and these block-face positions give the same full embedding
of $H$, including isolated components. All these checks take polynomial
time, for example $O(N^3)$ by direct searches. Hence unrestricted
$\SEFE$ belongs to NP. This is a membership proof, not a hardness result.
